% ===================== 1. GIỚI THIỆU =====================
\section{Giới thiệu}
\label{sec:intro}

\subsection{Bài toán và động lực}
\label{sec:motivation}
Bài toán xếp thời khóa biểu thuộc lớp \emph{phân bổ tài nguyên thời gian/phòng học}:
gán cho mỗi lớp một phòng và một tiết bắt đầu sao cho thỏa các ràng buộc tài nguyên
và tối đa hóa số lớp được xếp. Đây là một biến thể \emph{gán có giới hạn tài nguyên}
(resource-constrained assignment) --- NP-khó: các ràng buộc chống xung đột làm không
gian nghiệm khả thi vỡ vụn, và mỗi lớp có tới hàng trăm lựa chọn vị trí.

Động lực thực tế: cùng một bài toán nhưng xuất hiện ở nhiều quy mô và yêu cầu trái
ngược nhau --- từ xếp lịch một bộ môn (vài chục lớp, cần tối ưu tuyệt đối), tới
công cụ chỉnh lịch tương tác (cần phản hồi tức thời), tới xếp lịch toàn trường
(hàng nghìn lớp, chạy theo lô). Một phát hiện quan trọng dẫn dắt toàn bộ thiết kế:
\emph{nguồn độ khó không phải kích thước $N$ mà là độ chặt tài nguyên} --- điều
được kiểm chứng định lượng ở mục~\ref{sec:cpsat-limit}.

Hình~\ref{fig:problem} minh họa bài toán: gán mỗi lớp vào một khối $t_i$ tiết
\emph{liên tiếp} trong một phòng, tôn trọng ba ràng buộc cứng (sức chứa, không
trùng phòng, không trùng giáo viên) và không vắt chéo ngày.

\begin{figure}[H]
\centering
\begin{tikzpicture}[
  font=\scriptsize,
  gcell/.style={draw, gray!55, minimum size=0.5cm, inner sep=0pt},
  g1blk/.style={draw=cfast!70, fill=cfast!22, thick, rounded corners=1.5pt, minimum height=0.42cm, inner sep=1pt},
  g2blk/.style={draw=cexact!70, fill=cexact!18, thick, rounded corners=1.5pt, minimum height=0.42cm, inner sep=1pt},
  forb/.style={draw=cscale, fill=cscale!15, very thick, dashed, rounded corners=1.5pt, minimum height=0.42cm, inner sep=1pt, text=cscale},
  rlbl/.style={font=\scriptsize, anchor=east},
  note/.style={font=\scriptsize\itshape, align=center},
]
  % --- lưới phòng x tiết (một ngày: 12 tiết) ---
  \foreach \r in {0,1,2} \foreach \p in {1,...,12} {\node[gcell] at (\p*0.5,-\r*0.55){};}
  % nhãn tiết
  \foreach \p in {1,...,12} {\node[font=\tiny, anchor=south] at (\p*0.5,0.28){\p};}
  \node[font=\scriptsize\bfseries, anchor=south] at (3.25,0.62) {Thứ Hai --- 12 tiết};
  % nhãn phòng
  \node[rlbl] at (0.15,0) {$P_1\,(c{=}50)$};
  \node[rlbl] at (0.15,-0.55) {$P_2\,(c{=}80)$};
  \node[rlbl] at (0.15,-1.10) {$P_3\,(c{=}40)$};
  % --- các lớp đã xếp (màu = giáo viên) ---
  \node[g1blk, minimum width=0.95cm] (A) at (0.75,0)    {A\,$g_1$};
  \node[g2blk, minimum width=1.95cm] (B) at (1.75,-0.55){B\,$g_2$ ($s{=}70$)};
  \node[g1blk, minimum width=0.95cm] (C) at (2.25,-1.10){C\,$g_1$};
  % --- ranh giới ngày + khối vi phạm ---
  \draw[very thick, cscale!70] (6.25,0.42) -- (6.25,-1.45);
  \node[gcell, fill=gray!8] at (6.5,0){}; \node[gcell, fill=gray!8] at (7.0,0){};
  \node[forb, minimum width=1.45cm] (X) at (6.25,0) {\,$\times$\,};
  \node[note, text=cscale, anchor=west] at (6.7,0.0) {};
  % --- chú thích ràng buộc ---
  \node[note, anchor=north] (capn) at (1.75,-1.75) {Sức chứa: $s_i \le c_v$\\($B$ cần $P_2$, không vào được $P_3$)};
  \draw[->, thick, gray] (capn.north) to[out=90,in=270] (B.south);
  \node[note, text=cfast!70!black, anchor=north] (tch) at (5.0,-1.75) {Cùng giáo viên $g_1$:\\khác giờ (không trùng)};
  \draw[<->, thick, cfast, dashed] (A.east) to[out=-40,in=120] (C.north east);
  \node[note, text=cscale, anchor=north] (dayn) at (7.0,-1.0) {Không vắt\\chéo ngày};
  \draw[->, thick, cscale] (dayn.north) to[out=90,in=270] (X.south east);
\end{tikzpicture}
\caption{Minh họa bài toán trên một ngày (tuần có $5$ ngày $\times\,12$ tiết
$=60$ tiết). Mỗi lớp là một khối $t_i$ tiết liên tiếp đặt trong một phòng; màu khối
$=$ giáo viên. Ba ràng buộc cứng: \textbf{sức chứa} ($s_i\le c_v$), \textbf{không
trùng phòng} (mỗi ô tối đa một lớp), \textbf{không trùng giáo viên}; khối đỏ vi
phạm \emph{vắt chéo ngày}. Mục tiêu: tối đa hóa số lớp xếp được.}
\label{fig:problem}
\end{figure}

\subsection{Kiến trúc hệ thống (Đa tầng: Exact $\to$ Heuristic $\to$ Metaheuristic)}
\label{sec:arch}
Vì không một thuật toán nào tối ưu cho mọi quy mô, hệ thống được tổ chức thành
\emph{ba tầng} với một bộ định tuyến phân loại theo quy mô/độ chặt (Hình~\ref{fig:pipeline}).
Điểm thiết kế mấu chốt: cả ba tầng dùng chung nghiệm \textbf{Greedy} làm hạt giống
(warm-start cho CP-SAT, cấu hình nền cho Local Search và ALNS).

\begin{figure}[H]
\centering
\begin{tikzpicture}[
  font=\small,
  >={Stealth[round]},
  box/.style={rounded corners=3pt, draw, thick, align=center, inner sep=5pt, minimum height=1.0cm},
  io/.style={box, fill=soft, draw=accent},
  router/.style={box, fill=accent!12, draw=accent, minimum height=1.2cm},
  seed/.style={box, fill=cseed!12, draw=cseed, text=cseed!60!black},
  tier/.style={box, minimum width=3.7cm, text width=3.5cm},
  lbl/.style={font=\scriptsize\itshape, align=center},
  arr/.style={->, thick},
  seedarr/.style={->, thick, cseed, dashed},
]
\node[io, text width=1.7cm] (in) {Input\\\scriptsize $(N,M,$ lớp, phòng$)$};
\node[router, right=1.0cm of in, text width=2.2cm] (rt) {ROUTER\\\scriptsize định tuyến theo\\\scriptsize quy mô \& độ chặt};
\node[tier, fill=cexact!10, draw=cexact, text=cexact!50!black, right=2.0cm of rt, yshift=1.9cm] (t1)
  {\textbf{Exact}\\ CP-SAT\\[-1pt]\scriptsize tối ưu có chứng nhận};
\node[tier, fill=cfast!10, draw=cfast, text=cfast!55!black, right=2.0cm of rt] (t2)
  {\textbf{Heuristic}\\ Greedy $+$ Local Search\\[-1pt]\scriptsize mili-giây};
\node[tier, fill=cscale!10, draw=cscale, text=cscale!55!black, right=2.0cm of rt, yshift=-1.9cm] (t3)
  {\textbf{Metaheuristic}\\ Greedy $+$ ALNS\\[-1pt]\scriptsize chất lượng cao};
\node[io, right=1.8cm of t2, text width=1.7cm] (out) {Output\\\scriptsize lịch xếp $+\,Q$};
\node[seed, below=1.0cm of rt, text width=2.4cm] (gd) {Hạt giống \textbf{Greedy}\\\scriptsize (chung cho cả 3 tầng)};
\draw[arr] (in) -- (rt);
\draw[arr] (rt.east) -- (t1.west) node[lbl, pos=0.62, above, sloped] {nhỏ, cần tối ưu};
\draw[arr] (rt.east) -- (t2.west) node[lbl, pos=0.5, above] {tức thời};
\draw[arr] (rt.east) -- (t3.west) node[lbl, pos=0.62, below, sloped] {quy mô lớn};
\draw[arr] (t1.east) -- (out.west);
\draw[arr] (t2.east) -- (out.west);
\draw[arr] (t3.east) -- (out.west);
\draw[seedarr] (rt.south) -- (gd.north);
\draw[seedarr] (gd) to[out=35,in=255] (t1.south);
\draw[seedarr] (gd) to[out=10,in=260] (t2.south);
\draw[seedarr] (gd) to[out=-10,in=250] (t3.south);
\node[lbl, text=cseed, right=0.05cm of gd, yshift=-0.7cm] {\;mồi cả 3 tầng};
\end{tikzpicture}
\caption{Kiến trúc đa tầng \emph{Exact $\to$ Heuristic $\to$ Metaheuristic}. Bộ
\textsc{Router} định tuyến theo quy mô và độ chặt tài nguyên; nghiệm Greedy (nét
đứt tím) là hạt giống chung cho cả ba tầng.}
\label{fig:pipeline}
\end{figure}

\noindent Báo cáo đi theo đúng luồng dữ liệu của hệ thống: cấu trúc dữ liệu \&
trạng thái (Mục~\ref{sec:data}) $\to$ khởi tạo (Mục~\ref{sec:init}) $\to$ ba tầng
giải (Mục~\ref{sec:exact},~\ref{sec:ls},~\ref{sec:alns}) $\to$ tinh chỉnh tham số
(Mục~\ref{sec:tuning}) $\to$ thực nghiệm (Mục~\ref{sec:results}).

\subsection{Hàm mục tiêu và ràng buộc}
\label{sec:model}
Cấu trúc thời gian: $D=5$ ngày $\times\ H=12$ tiết, cho không gian
$\mathcal{T}=\{0,\dots,59\}$. Đầu vào: tập lớp $\mathcal{C}=\{1,\dots,N\}$ với mỗi
lớp $i$ có $(t_i,g_i,s_i)$ (thời lượng, giáo viên, sĩ số); tập phòng
$\mathcal{R}=\{1,\dots,M\}$ với sức chứa $c_v$.

\paragraph{Miền vị trí hợp lệ.} Vì lớp không vắt chéo ngày, tập tiết bắt đầu hợp lệ
của lớp thời lượng $t_i$ là
\begin{equation}
  \mathcal{S}(t_i) = \big\{\, d\cdot H + o \;\big|\; d\in\{0,\dots,D-1\},\;
  o\in\{0,\dots,H-t_i\} \,\big\}.
  \label{eq:allowed}
\end{equation}
Ràng buộc \emph{liên tục} và \emph{không xuyên ngày} được hấp thụ trọn vẹn vào miền
này --- không cần bất đẳng thức tường minh.

\paragraph{Biến quyết định, ràng buộc, mục tiêu.} Đặt $x_{i,r,u}=1$ nếu lớp $i$
xếp vào phòng $r$ bắt đầu tại $u\in\mathcal{S}(t_i)$, và $a_i=\sum_{r,u}x_{i,r,u}$.
\begin{align}
  &\textstyle\sum_{r}\sum_{u\in\mathcal{S}(t_i)} x_{i,r,u} = a_i \le 1
     && \forall i &&\text{(mỗi lớp tối đa một vị trí)}, \label{eq:c1}\\
  &x_{i,r,u}=0 \ \text{ nếu } c_r < s_i && && \text{(sức chứa)}, \label{eq:c2}\\
  &\textstyle\sum_{i}\sum_{u\le\tau<u+t_i} x_{i,r,u} \le 1
     && \forall r,\tau &&\text{(xung đột phòng)}, \label{eq:c3}\\
  &\textstyle\sum_{i:g_i=g}\sum_{r}\sum_{u\le\tau<u+t_i} x_{i,r,u} \le 1
     && \forall g,\tau &&\text{(xung đột giáo viên)}. \label{eq:c4}
\end{align}
\begin{equation}
  \boxed{\ \max\ Q = \sum_{i\in\mathcal{C}} a_i.\ }
  \label{eq:obj}
\end{equation}
